% Gliederung Stand 30.10.2026. Geplante Anhänge: A Messreihen, B Parameter/Topics, C Controller-Belegung, D Datenträger.
% Weitere Anhänge als eigene Dateien anlegen und in main.tex einbinden, sobald Inhalt da ist.
\chapter{Vollständige Messreihen}
\label{app:messreihen}

% ----------------------------------------------------------------------
% Abschnitt Testumgebung: am 25.11.2026 aus Kap. 5 (Abschnitt 5.3.3, sec:testumgebung) hierher verschoben (P6, Nils).
% Basis: G1@6a40e47 (Branch BA-Stable, Merge vom 23.11.2026), Dockerfiles und docker-compose.yml; Faktenblatt
% review/kap5-entwurf/fakten-testumgebung.md (d), dort noch Stand 25d5c1a. Docker-Dateien, unitree_mujoco/simulate_python
% und die Attrappen-Dateien sind zwischen 25d5c1a und 6a40e47 unverändert (git diff), nur test/run_tests.sh und neue Tests.
% Referenz-PC und Roboterstand liefert Nils (N4), bis dahin \pruefen.
% ----------------------------------------------------------------------
\section{Testumgebung der Versuche}
\label{anh:testumgebung}
Die Versuche aus Abschnitt~\ref{sec:versuchsplan} beziehen sich auf einen Referenz-PC, feste Softwarestände und einen Roboter, die Tabelle~\ref{tab:testumgebung} zusammenstellt.
Die Angaben zu den Softwareständen gelten für den Code-Stand der Arbeit.

% tab:testumgebung  (schlichter Stil, zwei Spalten). Art E.
\begin{table}[!ht]
    \setstretch{1.0}\tabstil\setlength{\tabcolsep}{4pt}
    \centering
    \caption[Testumgebung]{Testumgebung mit Referenz-PC, Softwareständen und Roboter. Angaben zu den Softwareständen gelten für den festgelegten Code-Stand.}
    \label{tab:testumgebung}
    \begin{tabularx}{\textwidth}{P{3.6cm}L}
        \toprule
        \tabkopf{Merkmal} & \tabkopf{Wert}\\
        \midrule
        Prozessor & \pruefen{Prozessor des Referenz-PCs von Nils}\\
        Arbeitsspeicher & \pruefen{Arbeitsspeicher des Referenz-PCs von Nils}\\
        Grafik & \pruefen{Grafikausstattung des Referenz-PCs von Nils}\\
        Betriebssystem & \pruefen{Betriebssystem und Version des Referenz-PCs von Nils}\\
        Container-Laufzeit & \pruefen{Version der Container-Laufzeit am Referenz-PC}\\
        % Basis: G1@6a40e47 g1pilot/docker/Dockerfile.sim:1, Dockerfile.real:7 (Basis-Image ros:humble ohne Prüfsumme)
        ROS-2-Version & ROS~2 Humble (Basis-Image über Tag, nicht auf eine Prüfsumme festgelegt). Paketstände nicht festgelegt \pruefen{Paketstände im Container auslesen und eintragen}\\
        % Basis: G1@6a40e47 g1pilot/docker/Dockerfile.mujoco:14 (pip install mujoco ohne Versionsangabe)
        MuJoCo-Version & nicht festgelegt \pruefen{installierte Version im Container auslesen und eintragen}\\
        % Basis: G1@6a40e47 g1pilot/docker/Dockerfile.sim:36-39, Dockerfile.real:42-45, Dockerfile.mujoco:16-19 (Commit 7c661d27f4ae)
        Version des Hersteller-SDK & Python-SDK in einer Abspaltung (Fork), Commit 7c661d2 (in allen drei Images gleich)\\
        % Basis: G1@6a40e47 g1pilot/docker/Dockerfile.sim:15,81, Dockerfile.real:21,80 (rmw_cyclonedds_cpp, apt ohne Version)
        DDS-Implementierung & Cyclone~DDS über \texttt{rmw\_cyclonedds\_cpp}, Version nicht festgelegt \pruefen{installierte Version im Container auslesen}\\
        % Basis: G1@6a40e47 g1pilot/docker/Dockerfile.sim:48 (Kommentar cp310 für Ubuntu 22.04 = ros:humble)
        Python & folgt dem Basis-Image, nicht festgelegt \pruefen{installierte Version im Container auslesen}\\
        % Basis: G1 git log -1 6a40e47 (Merge PR #59 vom 23.11.2026), Branch BA-Stable laut CLAUDE.md (Nils, 24.11.)
        Code-Stand der Umgebung & Repository \texttt{G1}, Branch \texttt{BA-Stable}, Commit 6a40e47 vom 23.11.2026\\
        Roboter (Modell, Ausführung, Softwarestand) & Unitree G1 \pruefen{Ausführung und Softwarestand des Roboters von Nils}\\
        Hände & \pruefen{Typ und Ausführung der Hände von Nils}\\
        \bottomrule
    \end{tabularx}
\end{table}

% Basis: fakten-testumgebung.md (d); G1@6a40e47 g1pilot/docker/Dockerfile.sim, Dockerfile.real, Dockerfile.mujoco, docker-compose.yml
Die Softwarestände legen die Container-Images der Umgebung fest (Abschnitt~\ref{sec:container}), die die Software bereits installiert, konfiguriert und getestet enthalten \cite{Boettiger2015}\ev{Boettiger2015.06}.
Die Versionen von MuJoCo und weiteren Python-Paketen sind darin nicht festgelegt und werden aus den gebauten Images ausgelesen \pruefen{installierte Versionen von MuJoCo, CycloneDDS und Pinocchio aus den gebauten Images ablesen und in tab:testumgebung eintragen}.
% Basis: G1 git diff --name-status 25d5c1a 6a40e47 (keine neuen Messskripte), todo-nils.md Kap. 5 „Festen Code-Stand im G1-Repo markieren“, N11
Die Messmittel für V4, V7 und V8 sind im Code-Stand 6a40e47 noch nicht enthalten, sodass diese Versuche an einem später markierten Stand laufen \pruefen{Tag bzw. Commit des Mess-Stands eintragen und die Testzahlen dort neu zählen}.

% Basis: G1@6a40e47, gezählt wie fakten-testumgebung.md (a), (c):
%   Dateien: git ls-tree -r --name-only 6a40e47 g1pilot/test/ | grep "/test_[^/]*\.py$" -> 55
%   Testfunktionen: grep -cE "^\s*(async )?def test_" je Datei, Summe -> 714 (keine Testklassen)
%   Gruppe sicherheit: Liste in g1pilot/test/run_tests.sh:21-26 (26 Dateien, alle vorhanden) -> 465 Testfunktionen
%   Dateien mit pytest.importorskip -> 38 (mit skipif zusätzlich: 39)
%   Zum Vergleich 25d5c1a mit derselben Methode: 46 / 597 / 25 Dateien, 442 / 30.
%   Ohne Roboter und Simulation: g1pilot/docs/60-A-Tests.md:21; still übersprungen, kein Nachweis: run_tests.sh:13-14, 60-A-Tests.md:59-61.
Am Stand 6a40e47 umfasst die Testsuite, die ohne Roboter und ohne laufende Simulation läuft, 714 Testfunktionen in 55 Dateien, davon die Gruppe „sicherheit“ 26 Dateien mit 465 Testfunktionen.
Fehlen Abhängigkeiten, überspringt die Suite Tests ohne Fehlermeldung, und 38 der 55 Dateien setzen solche Abhängigkeiten voraus.
Für K5.5 zählt deshalb nur ein Lauf ohne übersprungene Tests \pruefen{Testumgebung angeben (Image oder Python-Umgebung mit Version), in der die Suite für K5.5 läuft}.

% Basis: fakten-testumgebung.md (b); G1@6a40e47 g1pilot/test/test_arm_controller_real.py:29-160 (FakeG1, Stubs),
% g1pilot/test/test_loco_client_real.py:28-110 (FakeLoco), beide gegen 25d5c1a unverändert.
Die Attrappe aus Abschnitt~\ref{sec:testumgebung} bildet für die Arme den Zustand und die Armschnittstelle des Roboters nach, wobei die gemessenen Gelenkwinkel dem Sollwert mit einem einfachen Nachlauf folgen.
Für die Fortbewegung zeichnet sie die Aufrufe an den Onboard-Regler auf und führt dessen Betriebszustand als Zahl, ohne eine Bewegung zu simulieren.
Netzwerk, Beine und Gleichgewicht, Hände und Kollisionen bildet sie nicht nach, sodass Ergebnisse gegen sie die Logik der Programme belegen, nicht die Wirkung am Roboter.
% Basis: fakten-testumgebung.md (e) V5 (g1pilot/tools/walk_diag.py record, nur lesend, in 6a40e47 vorhanden); plan-kap5.md Abschnitt 7.3 und N11 (Messmittel baut Nils)
Für die Fortbewegung am Roboter steht eine nur lesende Aufzeichnung bereit, die übrigen Messmittel einschließlich der Zeitmessungen für V8 ergänzt die Arbeit \pruefen{Messmittel (N11) nach Fertigstellung bestätigen und Ablage für den Anhang nennen}.
